% Options for packages loaded elsewhere
% Options for packages loaded elsewhere
\PassOptionsToPackage{unicode}{hyperref}
\PassOptionsToPackage{hyphens}{url}
\PassOptionsToPackage{dvipsnames,svgnames,x11names}{xcolor}
%
\documentclass[
  letterpaper,
]{article}
\usepackage{xcolor}
\usepackage{amsmath,amssymb}
\setcounter{secnumdepth}{5}
\usepackage{iftex}
\ifPDFTeX
  \usepackage[T1]{fontenc}
  \usepackage[utf8]{inputenc}
  \usepackage{textcomp} % provide euro and other symbols
\else % if luatex or xetex
  \usepackage{unicode-math} % this also loads fontspec
  \defaultfontfeatures{Scale=MatchLowercase}
  \defaultfontfeatures[\rmfamily]{Ligatures=TeX,Scale=1}
\fi
\usepackage{lmodern}
\ifPDFTeX\else
  % xetex/luatex font selection
\fi
% Use upquote if available, for straight quotes in verbatim environments
\IfFileExists{upquote.sty}{\usepackage{upquote}}{}
\IfFileExists{microtype.sty}{% use microtype if available
  \usepackage[]{microtype}
  \UseMicrotypeSet[protrusion]{basicmath} % disable protrusion for tt fonts
}{}
\makeatletter
\@ifundefined{KOMAClassName}{% if non-KOMA class
  \IfFileExists{parskip.sty}{%
    \usepackage{parskip}
  }{% else
    \setlength{\parindent}{0pt}
    \setlength{\parskip}{6pt plus 2pt minus 1pt}}
}{% if KOMA class
  \KOMAoptions{parskip=half}}
\makeatother
% Make \paragraph and \subparagraph free-standing
\makeatletter
\ifx\paragraph\undefined\else
  \let\oldparagraph\paragraph
  \renewcommand{\paragraph}{
    \@ifstar
      \xxxParagraphStar
      \xxxParagraphNoStar
  }
  \newcommand{\xxxParagraphStar}[1]{\oldparagraph*{#1}\mbox{}}
  \newcommand{\xxxParagraphNoStar}[1]{\oldparagraph{#1}\mbox{}}
\fi
\ifx\subparagraph\undefined\else
  \let\oldsubparagraph\subparagraph
  \renewcommand{\subparagraph}{
    \@ifstar
      \xxxSubParagraphStar
      \xxxSubParagraphNoStar
  }
  \newcommand{\xxxSubParagraphStar}[1]{\oldsubparagraph*{#1}\mbox{}}
  \newcommand{\xxxSubParagraphNoStar}[1]{\oldsubparagraph{#1}\mbox{}}
\fi
\makeatother

\usepackage{color}
\usepackage{fancyvrb}
\newcommand{\VerbBar}{|}
\newcommand{\VERB}{\Verb[commandchars=\\\{\}]}
\DefineVerbatimEnvironment{Highlighting}{Verbatim}{commandchars=\\\{\}}
% Add ',fontsize=\small' for more characters per line
\usepackage{framed}
\definecolor{shadecolor}{RGB}{241,243,245}
\newenvironment{Shaded}{\begin{snugshade}}{\end{snugshade}}
\newcommand{\AlertTok}[1]{\textcolor[rgb]{0.68,0.00,0.00}{#1}}
\newcommand{\AnnotationTok}[1]{\textcolor[rgb]{0.37,0.37,0.37}{#1}}
\newcommand{\AttributeTok}[1]{\textcolor[rgb]{0.40,0.45,0.13}{#1}}
\newcommand{\BaseNTok}[1]{\textcolor[rgb]{0.68,0.00,0.00}{#1}}
\newcommand{\BuiltInTok}[1]{\textcolor[rgb]{0.00,0.23,0.31}{#1}}
\newcommand{\CharTok}[1]{\textcolor[rgb]{0.13,0.47,0.30}{#1}}
\newcommand{\CommentTok}[1]{\textcolor[rgb]{0.37,0.37,0.37}{#1}}
\newcommand{\CommentVarTok}[1]{\textcolor[rgb]{0.37,0.37,0.37}{\textit{#1}}}
\newcommand{\ConstantTok}[1]{\textcolor[rgb]{0.56,0.35,0.01}{#1}}
\newcommand{\ControlFlowTok}[1]{\textcolor[rgb]{0.00,0.23,0.31}{\textbf{#1}}}
\newcommand{\DataTypeTok}[1]{\textcolor[rgb]{0.68,0.00,0.00}{#1}}
\newcommand{\DecValTok}[1]{\textcolor[rgb]{0.68,0.00,0.00}{#1}}
\newcommand{\DocumentationTok}[1]{\textcolor[rgb]{0.37,0.37,0.37}{\textit{#1}}}
\newcommand{\ErrorTok}[1]{\textcolor[rgb]{0.68,0.00,0.00}{#1}}
\newcommand{\ExtensionTok}[1]{\textcolor[rgb]{0.00,0.23,0.31}{#1}}
\newcommand{\FloatTok}[1]{\textcolor[rgb]{0.68,0.00,0.00}{#1}}
\newcommand{\FunctionTok}[1]{\textcolor[rgb]{0.28,0.35,0.67}{#1}}
\newcommand{\ImportTok}[1]{\textcolor[rgb]{0.00,0.46,0.62}{#1}}
\newcommand{\InformationTok}[1]{\textcolor[rgb]{0.37,0.37,0.37}{#1}}
\newcommand{\KeywordTok}[1]{\textcolor[rgb]{0.00,0.23,0.31}{\textbf{#1}}}
\newcommand{\NormalTok}[1]{\textcolor[rgb]{0.00,0.23,0.31}{#1}}
\newcommand{\OperatorTok}[1]{\textcolor[rgb]{0.37,0.37,0.37}{#1}}
\newcommand{\OtherTok}[1]{\textcolor[rgb]{0.00,0.23,0.31}{#1}}
\newcommand{\PreprocessorTok}[1]{\textcolor[rgb]{0.68,0.00,0.00}{#1}}
\newcommand{\RegionMarkerTok}[1]{\textcolor[rgb]{0.00,0.23,0.31}{#1}}
\newcommand{\SpecialCharTok}[1]{\textcolor[rgb]{0.37,0.37,0.37}{#1}}
\newcommand{\SpecialStringTok}[1]{\textcolor[rgb]{0.13,0.47,0.30}{#1}}
\newcommand{\StringTok}[1]{\textcolor[rgb]{0.13,0.47,0.30}{#1}}
\newcommand{\VariableTok}[1]{\textcolor[rgb]{0.07,0.07,0.07}{#1}}
\newcommand{\VerbatimStringTok}[1]{\textcolor[rgb]{0.13,0.47,0.30}{#1}}
\newcommand{\WarningTok}[1]{\textcolor[rgb]{0.37,0.37,0.37}{\textit{#1}}}

\usepackage{longtable,booktabs,array}
\newcounter{none} % for unnumbered tables
\usepackage{calc} % for calculating minipage widths
% Correct order of tables after \paragraph or \subparagraph
\usepackage{etoolbox}
\makeatletter
\patchcmd\longtable{\par}{\if@noskipsec\mbox{}\fi\par}{}{}
\makeatother
% Allow footnotes in longtable head/foot
\IfFileExists{footnotehyper.sty}{\usepackage{footnotehyper}}{\usepackage{footnote}}
\makesavenoteenv{longtable}
\usepackage{graphicx}
\makeatletter
\newsavebox\pandoc@box
\newcommand*\pandocbounded[1]{% scales image to fit in text height/width
  \sbox\pandoc@box{#1}%
  \Gscale@div\@tempa{\textheight}{\dimexpr\ht\pandoc@box+\dp\pandoc@box\relax}%
  \Gscale@div\@tempb{\linewidth}{\wd\pandoc@box}%
  \ifdim\@tempb\p@<\@tempa\p@\let\@tempa\@tempb\fi% select the smaller of both
  \ifdim\@tempa\p@<\p@\scalebox{\@tempa}{\usebox\pandoc@box}%
  \else\usebox{\pandoc@box}%
  \fi%
}
% Set default figure placement to htbp
\def\fps@figure{htbp}
\makeatother


% definitions for citeproc citations
\NewDocumentCommand\citeproctext{}{}
\NewDocumentCommand\citeproc{mm}{%
  \begingroup\def\citeproctext{#2}\cite{#1}\endgroup}
\makeatletter
 % allow citations to break across lines
 \let\@cite@ofmt\@firstofone
 % avoid brackets around text for \cite:
 \def\@biblabel#1{}
 \def\@cite#1#2{{#1\if@tempswa , #2\fi}}
\makeatother
\newlength{\cslhangindent}
\setlength{\cslhangindent}{1.5em}
\newlength{\csllabelwidth}
\setlength{\csllabelwidth}{3em}
\newenvironment{CSLReferences}[2] % #1 hanging-indent, #2 entry-spacing
 {\begin{list}{}{%
  \setlength{\itemindent}{0pt}
  \setlength{\leftmargin}{0pt}
  \setlength{\parsep}{0pt}
  % turn on hanging indent if param 1 is 1
  \ifodd #1
   \setlength{\leftmargin}{\cslhangindent}
   \setlength{\itemindent}{-1\cslhangindent}
  \fi
  % set entry spacing
  \setlength{\itemsep}{#2\baselineskip}}}
 {\end{list}}
\usepackage{calc}
\newcommand{\CSLBlock}[1]{\hfill\break\parbox[t]{\linewidth}{\strut\ignorespaces#1\strut}}
\newcommand{\CSLLeftMargin}[1]{\parbox[t]{\csllabelwidth}{\strut#1\strut}}
\newcommand{\CSLRightInline}[1]{\parbox[t]{\linewidth - \csllabelwidth}{\strut#1\strut}}
\newcommand{\CSLIndent}[1]{\hspace{\cslhangindent}#1}



\setlength{\emergencystretch}{3em} % prevent overfull lines

\providecommand{\tightlist}{%
  \setlength{\itemsep}{0pt}\setlength{\parskip}{0pt}}



 


\usepackage{orcidlink}  
\usepackage[nonatbib, preprint]{neurips_2023}
\usepackage{amsmath}
\usepackage{amsthm}
\usepackage[utf8]{inputenc} % allow utf-8 input
\usepackage[T1]{fontenc}    % use 8-bit T1 fonts
\usepackage{hyperref}       % hyperlinks
\usepackage{url}            % simple URL typesetting
\usepackage{booktabs}       % professional-quality tables
\usepackage{amsfonts}       % blackboard math symbols
\usepackage{nicefrac}       % compact symbols for 1/2, etc.
\usepackage{microtype}      % microtypography
\usepackage{xcolor}         % colors
\usepackage[numbers]{natbib}
\usepackage{comment}
\usepackage{graphicx} 
\bibliographystyle{abbrvnat}
\makeatletter
\@ifpackageloaded{caption}{}{\usepackage{caption}}
\AtBeginDocument{%
\ifdefined\contentsname
  \renewcommand*\contentsname{Table of contents}
\else
  \newcommand\contentsname{Table of contents}
\fi
\ifdefined\listfigurename
  \renewcommand*\listfigurename{List of Figures}
\else
  \newcommand\listfigurename{List of Figures}
\fi
\ifdefined\listtablename
  \renewcommand*\listtablename{List of Tables}
\else
  \newcommand\listtablename{List of Tables}
\fi
\ifdefined\figurename
  \renewcommand*\figurename{Figure}
\else
  \newcommand\figurename{Figure}
\fi
\ifdefined\tablename
  \renewcommand*\tablename{Table}
\else
  \newcommand\tablename{Table}
\fi
}
\@ifpackageloaded{float}{}{\usepackage{float}}
\floatstyle{ruled}
\@ifundefined{c@chapter}{\newfloat{codelisting}{h}{lop}}{\newfloat{codelisting}{h}{lop}[chapter]}
\floatname{codelisting}{Listing}
\newcommand*\listoflistings{\listof{codelisting}{List of Listings}}
\makeatother
\makeatletter
\makeatother
\makeatletter
\@ifpackageloaded{caption}{}{\usepackage{caption}}
\@ifpackageloaded{subcaption}{}{\usepackage{subcaption}}
\makeatother
\makeatletter
\@ifpackageloaded{algorithm}{}{\usepackage{algorithm}}
\makeatother
\makeatletter
\@ifpackageloaded{algpseudocode}{}{\usepackage{algpseudocode}}
\makeatother
\usepackage{bookmark}
\IfFileExists{xurl.sty}{\usepackage{xurl}}{} % add URL line breaks if available
\urlstyle{same}
\hypersetup{
  pdftitle={Syntax-Guided Program Reduction: A Survey and the unfluff Implementation},
  pdfauthor={Unfluff Contributors},
  pdfkeywords={program reduction, test case
minimization, syntax-guided, Perses, delta
debugging, debugging, survey},
  colorlinks=true,
  linkcolor={blue},
  filecolor={Maroon},
  citecolor={Blue},
  urlcolor={Blue},
  pdfcreator={LaTeX via pandoc}}



\title{Syntax-Guided Program Reduction: A Survey and the unfluff
Implementation}
\author{
\textbf{Unfluff Contributors}%
\\%
%
    %
    Theseus Ship\\%
    %
%
}
\date{2026-06-13}
\begin{document}
\maketitle
\begin{abstract}
Program reduction---the automatic minimization of a failing test case to
the smallest input that preserves the failure---is a critical tool in
software debugging, compiler testing, and fuzzing. Given a program \(P\)
that exhibits a property \(\psi\) (e.g., triggers a compiler crash),
program reduction seeks the smallest \(P'\) such that \(\psi(P')\)
holds. This paper surveys the landscape of program reduction tools and
techniques, from the foundational delta debugging algorithm through
syntax-guided approaches like Perses and token-level methods like T-Rec.
We present \textbf{unfluff}, a high-performance Rust implementation of
the Perses algorithm that leverages tree-sitter for multi-language
support across seven languages (Python, JavaScript, TypeScript, Rust,
Go, C, C++). unfluff provides both a Rust-based primary implementation
achieving 3--13x speedups over Python and a Python reference
implementation for accessibility. We evaluate unfluff against
established reducers, discuss its design decisions and their rationale,
and identify open challenges and future research directions in program
reduction.
\end{abstract}

\floatname{algorithm}{Algorithm}


\section{1. Introduction}\label{introduction}

\subsection{1.1 The Problem}\label{the-problem}

Program reduction is important and widely demanded. Given a program
\(P\) that exhibits a property \(\psi\)---for example, a C program that
crashes GCC when it is being compiled---the objective of program
reduction is to generate a smaller program \(P'\) from \(P\) that still
exhibits the same property. The difficulty is combinatorial: a 100-line
file has \(2^{100}\) subsets. Tools must exploit structure (syntax,
grammar, program semantics) to navigate this space efficiently.

The ideal is global minimality: finding the absolute smallest program
that preserves the property. However, global minimality is
computationally intractable in general, so practical tools target weaker
guarantees:

{\def\LTcaptype{none} % do not increment counter
\begin{longtable}[]{@{}
  >{\raggedright\arraybackslash}p{(\linewidth - 2\tabcolsep) * \real{0.1895}}
  >{\raggedright\arraybackslash}p{(\linewidth - 2\tabcolsep) * \real{0.8105}}@{}}
\toprule\noalign{}
\begin{minipage}[b]{\linewidth}\raggedright
Minimality Level
\end{minipage} & \begin{minipage}[b]{\linewidth}\raggedright
Definition
\end{minipage} \\
\midrule\noalign{}
\endhead
\bottomrule\noalign{}
\endlastfoot
\textbf{1-minimal} & No single element can be removed while preserving
\(\psi\) \\
\textbf{1-tree-minimal} & No AST node can be further simplified (delete
or unwrap) while preserving \(\psi\) \\
\textbf{Global minimum} & No smaller program preserves \(\psi\)
(computationally intractable in general) \\
\end{longtable}
}

\subsection{1.2 Applications}\label{applications}

Program reduction serves several critical applications:

\begin{enumerate}
\def\labelenumi{\arabic{enumi}.}
\item
  \textbf{Compiler bug reporting:} When a compiler crashes or
  miscompiles, reducing the trigger to a minimal test case helps
  developers understand and fix the bug. The GCC and LLVM projects both
  maintain guidelines for submitting reduced test cases.
\item
  \textbf{Fuzzing:} Grammar-based fuzzers generate many inputs that
  trigger bugs. Reducing these inputs to minimal triggers helps
  developers understand the root cause.
\item
  \textbf{Regression testing:} Reduced test cases are easier to
  maintain, faster to execute, and clearer in their intent.
\item
  \textbf{Language specification:} Reduced examples help language
  designers understand edge cases and ambiguities in language
  specifications.
\end{enumerate}

\subsection{1.3 Contributions}\label{contributions}

This paper makes the following contributions:

\begin{enumerate}
\def\labelenumi{\arabic{enumi}.}
\tightlist
\item
  \textbf{A comprehensive survey} of program reduction tools and
  techniques, from delta debugging through syntax-guided approaches to
  token-level methods.
\item
  \textbf{The unfluff system}, a high-performance Rust implementation of
  the Perses algorithm with multi-language support via tree-sitter.
\item
  \textbf{An evaluation} of unfluff against established reducers,
  including performance benchmarks and reduction quality comparisons.
\item
  \textbf{An analysis} of design decisions, trade-offs, and open
  challenges in program reduction.
\end{enumerate}

\subsection{1.4 Paper Organization}\label{paper-organization}

The remainder of this paper is organized as follows. Section 2 surveys
background and related work. Section 3 describes the unfluff system
architecture and implementation. Section 4 evaluates unfluff against
established tools. Section 5 discusses limitations and future work.
Section 6 concludes.

\section{2. Background and Related
Work}\label{background-and-related-work}

\subsection{2.1 Delta Debugging (ddmin)}\label{delta-debugging-ddmin}

The foundational approach by Zeller and Hildebrandt (2002). Given a list
of elements (lines, tokens, characters), ddmin repeatedly halves
partitions and tests whether each half preserves the property.

\textbf{Algorithm:} The ddmin algorithm operates as follows:

\begin{enumerate}
\def\labelenumi{\arabic{enumi}.}
\tightlist
\item
  \textbf{Split Phase:} Split the input \(L\) into \(n\) partitions. For
  each partition \(u\), test if \(u\) alone preserves \(\psi\). If yes,
  remove the complement.
\item
  \textbf{Complement Phase:} Test if the complement of each partition
  preserves \(\psi\). If yes, remove \(u\).
\item
  \textbf{Refine:} Try to split each remaining partition into halves
  (\(n \to 2n\)). Resume at Step 1.
\item
  \textbf{Terminate:} When each partition cannot be further split, the
  remaining elements form the reduced result.
\end{enumerate}

\textbf{Properties:}

\begin{itemize}
\tightlist
\item
  Guarantees 1-minimality.
\item
  \(O(n^2)\) test invocations in the worst case.
\item
  Language-independent: operates on flat sequences.
\end{itemize}

\textbf{Weaknesses:}

\begin{itemize}
\tightlist
\item
  Produces syntactically invalid intermediates (wasted test invocations
  on broken code).
\item
  No structural awareness: cannot distinguish between removing a keyword
  vs.~removing an entire statement.
\item
  High test overhead: each test requires invoking the program under
  test.
\end{itemize}

\subsection{2.2 Hierarchical Delta Debugging
(HDD)}\label{hierarchical-delta-debugging-hdd}

Misherghi and Su (2006) applies ddmin hierarchically: first at the line
level, then at the AST level, then at finer granularities.

\textbf{Algorithm:} HDD operates in multiple passes:

\begin{enumerate}
\def\labelenumi{\arabic{enumi}.}
\tightlist
\item
  \textbf{Line-level pass:} Apply ddmin to lines of the source file.
\item
  \textbf{Block-level pass:} Group lines into blocks (e.g., function
  bodies, compound statements) and apply ddmin to blocks.
\item
  \textbf{AST-level pass:} Parse the reduced file and apply ddmin to AST
  subtrees.
\end{enumerate}

\textbf{Properties:}

\begin{itemize}
\tightlist
\item
  Exploits parse tree structure to reduce the search space.
\item
  Better results than plain ddmin: 2--10x smaller outputs on average.
\end{itemize}

\textbf{Weaknesses:}

\begin{itemize}
\tightlist
\item
  Still produces invalid intermediates at each level.
\item
  Requires a full parser (language-specific dependency).
\item
  The hierarchy of levels is ad-hoc: no formal basis for choosing
  granularity.
\end{itemize}

\subsection{2.3 C-Reduce}\label{c-reduce}

Regehr et al. (2012) is the de facto standard for C/C++ reduction. A
Perl-based tool applying approximately 80 specialized transformation
passes via \texttt{clang\_delta}, \texttt{unifdef}, \texttt{clex}, and
others.

\textbf{Architecture:}

\begin{itemize}
\tightlist
\item
  \textbf{Transformation passes:} Each pass applies a specific reduction
  (e.g., removing a function parameter, simplifying a switch statement,
  inlining a function).
\item
  \textbf{Interestingness test:} The user provides a script that exits 0
  when the bug is still present.
\item
  \textbf{Parallel execution:} Multiple candidates are tested
  concurrently using fork/eval.
\end{itemize}

\textbf{Properties:}

\begin{itemize}
\tightlist
\item
  Extremely effective for C/C++: widely used in compiler bug reporting.
\item
  Rich set of language-specific transformations.
\item
  Parallel execution provides near-linear speedup.
\end{itemize}

\textbf{Weaknesses:}

\begin{itemize}
\tightlist
\item
  Heavily C/C++-centric: works ``pretty well'' on other languages but is
  not optimized for them.
\item
  Complex Perl codebase: 2,071 commits of accumulated complexity.
\item
  No formal guarantee of syntactic validity at every step.
\item
  Requires Clang as a dependency.
\end{itemize}

\subsection{2.4 Shrinkray}\label{shrinkray}

DRMacIver (2023) is a modern Python-based multiformat reducer by
DRMacIver. Generic algorithm with format-specific passes (C, Python,
JSON, DIMACS CNF).

\textbf{Architecture:}

\begin{itemize}
\tightlist
\item
  \textbf{Generic reduction:} Applies delta debugging to the entire
  input.
\item
  \textbf{Format-specific passes:} Language-specific transformations
  (e.g., Python AST simplification, C type system optimization).
\item
  \textbf{High parallelism:} Designed for concurrent test execution.
\item
  \textbf{Live UI:} Interactive display of reduction progress.
\end{itemize}

\textbf{Properties:}

\begin{itemize}
\tightlist
\item
  Multiformat: supports C, Python, JSON, DIMACS CNF.
\item
  Highly parallel: uses Python's asyncio for concurrent testing.
\item
  Modern Python codebase: clean architecture, well-documented.
\item
  Active development: 590 commits, ongoing maintenance.
\end{itemize}

\textbf{Weaknesses:}

\begin{itemize}
\tightlist
\item
  Generic algorithm: not syntax-guided in the Perses sense.
\item
  Format-specific passes are basic compared to C-Reduce.
\item
  No formal syntactic validity guarantee.
\end{itemize}

\subsection{2.5 Perses}\label{perses}

Sun et al. (2018) introduced a syntax-guided framework that exploits the
formal grammar of the language under reduction. Key innovations:

\subsubsection{2.5.1 Grammar Normalization to Perses Normal Form
(PNF)}\label{grammar-normalization-to-perses-normal-form-pnf}

Perses normalizes grammars to a form that distinguishes between:

\begin{itemize}
\tightlist
\item
  \textbf{Regular rules:} \(A \to B_1 B_2 \ldots B_n\) (fixed number of
  children).
\item
  \textbf{Kleene star:} \(A \to B_1^*\) (zero or more children).
\item
  \textbf{Kleene plus:} \(A \to B_1^+\) (one or more children).
\item
  \textbf{Optional:} \(A \to B_1?\) (zero or one child).
\end{itemize}

This normalization enables node-type-aware dispatch: Kleene nodes use
ddmin on children (efficient for list-like constructs), while regular
nodes use BoundedBFS with the subsume relation.

\subsubsection{2.5.2 Node-Type-Aware
Dispatch}\label{node-type-aware-dispatch}

The main reduction algorithm processes nodes in priority order (largest
first):

\begin{itemize}
\tightlist
\item
  \textbf{Kleene-star/optional nodes:} Use ddmin to delete children.
  Each child is independent in terms of syntax validity, so ddmin is
  efficient.
\item
  \textbf{Kleene-plus nodes:} Use ddmin with the constraint that at
  least one child must remain.
\item
  \textbf{Regular rule nodes:} Use BoundedBFS to find replacement
  candidates that are type-compatible with the target node.
\end{itemize}

\subsubsection{2.5.3 Reparse as Validity
Gate}\label{reparse-as-validity-gate}

Every transform is validated by reparsing the modified source. Only
transforms that produce a syntactically valid program (no ERROR/MISSING
nodes) proceed to interestingness testing. This ensures
\(P_{\text{invalid}} = \emptyset\) in the search space.

\subsubsection{2.5.4 Subsume Relation}\label{subsume-relation}

The subsume relation determines which nodes can replace a target node:

\[
B <: A \text{ if } B \text{ can be derived from } A
\]

For example, \texttt{if\_stmt\ \textless{}:\ stmt} because an
\texttt{if\_stmt} can appear wherever a \texttt{stmt} is expected. This
enables finding valid replacement candidates that are not exact type
matches.

\subsubsection{2.5.5 Evaluation}\label{evaluation}

Perses's results are respectively 2\% and 45\% in size of those from DD
and HDD, using 23\% and 47\% of the time. Compared to C-Reduce, Perses
takes only 38--60\% of the reduction time.

\subsection{2.6 Vulcan}\label{vulcan}

Xu et al. (2023) addresses a fundamental limitation of Perses:
1-tree-minimality is not global minimality. Vulcan introduces
transformations that can escape local minima that Perses cannot.

\textbf{Key Insight:} Some reductions require ``horizontal'' moves
(e.g., renaming a variable to a shorter name, simplifying an expression)
that are not strictly ``smaller'' but enable further reduction.

\textbf{Evaluation:} Vulcan produces results 13--61\% smaller than
Perses on some benchmarks, particularly for programs with many redundant
identifiers or complex expressions.

\subsection{2.7 T-Rec}\label{t-rec}

Xu et al. (2025) operates at the lexical token level rather than the AST
level. Exploits the observation that many reductions are ``horizontal''
(removing tokens within a node) rather than ``vertical'' (removing
entire nodes).

\textbf{Key Insight:} After tree-level reduction reaches a fixpoint, a
token-level pass can further reduce within nodes (e.g., removing unused
import names, shortening identifiers, removing redundant whitespace).

\textbf{Evaluation:} T-Rec complements tree-level tools and can make
reductions that Perses misses, particularly for programs with verbose
syntax.

\subsection{2.8 bonsai}\label{bonsai}

{nnunley} (2024) is a Rust reference implementation of Perses adapted
for tree-sitter grammars. Supports Python, JavaScript, and Rust. Uses
tree-sitter's supertype/subtype system and \texttt{locals.scm} queries
for scope-aware transforms.

\textbf{Key Features:}

\begin{itemize}
\tightlist
\item
  \textbf{Rust performance:} Native implementation for speed.
\item
  \textbf{tree-sitter integration:} Incremental parsing, 300+ language
  grammars.
\item
  \textbf{Scope-aware transforms:} Identifier unification, dead
  definition removal.
\item
  \textbf{Parallel execution:} Multiple candidates tested concurrently.
\end{itemize}

\subsection{2.9 Comparison Matrix}\label{comparison-matrix}

{\def\LTcaptype{none} % do not increment counter
\begin{longtable}[]{@{}
  >{\raggedright\arraybackslash}p{(\linewidth - 10\tabcolsep) * \real{0.1509}}
  >{\raggedright\arraybackslash}p{(\linewidth - 10\tabcolsep) * \real{0.2547}}
  >{\raggedright\arraybackslash}p{(\linewidth - 10\tabcolsep) * \real{0.1509}}
  >{\raggedright\arraybackslash}p{(\linewidth - 10\tabcolsep) * \real{0.0755}}
  >{\raggedright\arraybackslash}p{(\linewidth - 10\tabcolsep) * \real{0.1509}}
  >{\raggedright\arraybackslash}p{(\linewidth - 10\tabcolsep) * \real{0.2170}}@{}}
\toprule\noalign{}
\begin{minipage}[b]{\linewidth}\raggedright
Tool
\end{minipage} & \begin{minipage}[b]{\linewidth}\raggedright
Algorithm
\end{minipage} & \begin{minipage}[b]{\linewidth}\raggedright
Syntax-Guarantee
\end{minipage} & \begin{minipage}[b]{\linewidth}\raggedright
Parallel
\end{minipage} & \begin{minipage}[b]{\linewidth}\raggedright
Languages
\end{minipage} & \begin{minipage}[b]{\linewidth}\raggedright
Minimality
\end{minipage} \\
\midrule\noalign{}
\endhead
\bottomrule\noalign{}
\endlastfoot
\textbf{ddmin} & Binary search on partitions & None & Optional & Any &
1-minimal \\
\textbf{HDD} & Hierarchical ddmin & None & Optional & Any & 1-minimal \\
\textbf{C-Reduce} & 80+ transformation passes & Partial & Yes & C/C++ (+
others) & Beyond 1-minimal \\
\textbf{Shrinkray} & Generic + format passes & None & Yes & Multiformat
& 1-minimal \\
\textbf{Perses} & Priority queue + BoundedBFS & Yes & No & Any (grammar)
& 1-tree-minimal \\
\textbf{bonsai} & Perses via tree-sitter & Yes & Yes & Python, JS, Rust
& 1-tree-minimal \\
\textbf{Vulcan} & Perses + escape transforms & Yes & No & Any (grammar)
& Beyond 1-tree-minimal \\
\textbf{T-Rec} & Token-level Perses & Yes & No & Any & Beyond
1-tree-minimal \\
\textbf{unfluff} & Simplified Perses & Yes & No & 7 languages &
1-tree-minimal (simplified) \\
\end{longtable}
}

\section{3. The unfluff System}\label{the-unfluff-system}

\subsection{3.1 Design Philosophy}\label{design-philosophy}

unfluff is designed around three core principles:

\begin{enumerate}
\def\labelenumi{\arabic{enumi}.}
\tightlist
\item
  \textbf{Accessibility:} Minimal dependencies, Python-native
  interfaces, and clear documentation lower the barrier to adoption.
\item
  \textbf{Performance:} A Rust-based primary implementation provides the
  speed needed for large-scale use.
\item
  \textbf{Extensibility:} A modular architecture makes it
  straightforward to add new languages, transforms, or interestingness
  test interfaces.
\end{enumerate}

\subsection{3.2 Architecture Overview}\label{architecture-overview}

unfluff implements the Perses algorithm with the following data flow:

\begin{verbatim}
Input File → Parse → CST → Generate Candidates → Apply Transform
         → Reparse+Validate → Interestingness Test → Accept/Reject → Loop
\end{verbatim}

The system consists of three main subsystems:

\subsubsection{3.2.1 Parsing Subsystem}\label{parsing-subsystem}

\begin{itemize}
\tightlist
\item
  \textbf{\texttt{parser.rs}:} Parses source files into concrete syntax
  trees (CSTs) using tree-sitter. Walks the tree to extract node
  information (kind, byte range, token count, error status, child
  information).
\item
  \textbf{\texttt{grammar.rs}:} Manages tree-sitter language loading,
  type queries (supertypes, subtypes), and Kleene node detection.
\end{itemize}

\subsubsection{3.2.2 Reduction Subsystem}\label{reduction-subsystem}

\begin{itemize}
\tightlist
\item
  \textbf{\texttt{reducer.rs}:} Implements the priority-queue-based
  reduction loop. Manages the worklist, dispatches to transform
  functions, and handles interestingness testing.
\item
  \textbf{\texttt{transforms.rs}:} Implements Delete, Unwrap, and
  BoundedBFS transforms. Generates candidates and applies transforms
  with reparse validation.
\end{itemize}

\subsubsection{3.2.3 Token Reduction
Subsystem}\label{token-reduction-subsystem}

\begin{itemize}
\tightlist
\item
  \textbf{\texttt{token\_reduce.rs}:} Applies fine-grained
  transformations after tree-level reduction reaches a fixpoint:
  trailing whitespace removal, redundant parenthesis removal,
  unnecessary semicolon removal.
\end{itemize}

\subsubsection{3.2.4 Supporting Modules}\label{supporting-modules}

\begin{itemize}
\tightlist
\item
  \textbf{\texttt{cache.rs}:} blake3-based content-hash caching to avoid
  re-running identical interestingness tests.
\item
  \textbf{\texttt{checker.rs}:} Static analysis checks (RED100--RED204
  rules) for identifying reducible patterns.
\item
  \textbf{\texttt{shrink.rs}:} Shrinkray-compatible interface for
  integration with existing Shrinkray test scripts.
\item
  \textbf{\texttt{tree.rs}:} Core data types (\texttt{NodeInfo},
  \texttt{TransformCandidate}, \texttt{TransformKind},
  \texttt{ParseResult}).
\end{itemize}

\subsection{3.3 Algorithm in Detail}\label{algorithm-in-detail}

\subsubsection{3.3.1 Candidate Generation}\label{candidate-generation}

For each node in the CST (excluding the root and error nodes):

\begin{enumerate}
\def\labelenumi{\arabic{enumi}.}
\tightlist
\item
  \textbf{Check if Kleene node:} If all named children are of the same
  type, the node is a Kleene container. Generate a ddmin candidate.
\item
  \textbf{Otherwise:} Generate a Delete candidate.
\end{enumerate}

\subsubsection{3.3.2 Priority Queue
Ordering}\label{priority-queue-ordering}

Following Perses, candidates are ordered by token count (largest
subtrees first). This strategy maximizes the potential reduction per
test invocation:

\[
\text{priority}(c) = \text{token\_count}(c.\text{target})
\]

\subsubsection{3.3.3 Transform Application}\label{transform-application}

For each candidate in priority order:

\begin{enumerate}
\def\labelenumi{\arabic{enumi}.}
\tightlist
\item
  \textbf{Delete:} Remove the target node's bytes from the source.
\item
  \textbf{Unwrap:} Replace the target node with a type-compatible
  descendant (discovered via BoundedBFS).
\item
  \textbf{ddmin:} Apply delta debugging to the children of a Kleene
  node.
\end{enumerate}

\subsubsection{3.3.4 Reparse Validation}\label{reparse-validation}

After applying a transform, reparse the modified source:

\[
\text{result} = \text{parse}(\text{modified\_source})
\]

Accept only if \(E' \leq E_0\) (no new errors introduced), where \(E_0\)
is the initial error count.

\subsubsection{3.3.5 Interestingness
Testing}\label{interestingness-testing}

If the transform passes reparse validation, test whether the property is
preserved:

\[
\text{is\_interesting} = \psi(\text{modified\_source})
\]

If true, accept the transform and add the target node's children to the
worklist. Otherwise, reject and try the next candidate.

\subsubsection{3.3.6 Fixpoint Iteration}\label{fixpoint-iteration}

Repeat until no more transforms are accepted. The final result is
1-tree-minimal (no single node can be further simplified while
preserving the property).

\subsection{3.4 Transforms in Detail}\label{transforms-in-detail}

\subsubsection{3.4.1 Delete Transform}\label{delete-transform}

The simplest transform: remove an entire node from the AST. Applicable
to nodes that are not required by their parent's grammar rule.

\textbf{Example:} Given a function definition with a docstring:

\begin{Shaded}
\begin{Highlighting}[]
\KeywordTok{def}\NormalTok{ f(x):}
    \CommentTok{"""Compute something."""}
    \ControlFlowTok{return}\NormalTok{ x }\OperatorTok{+} \DecValTok{1}
\end{Highlighting}
\end{Shaded}

Deleting the docstring node produces:

\begin{Shaded}
\begin{Highlighting}[]
\KeywordTok{def}\NormalTok{ f(x):}
    \ControlFlowTok{return}\NormalTok{ x }\OperatorTok{+} \DecValTok{1}
\end{Highlighting}
\end{Shaded}

\subsubsection{3.4.2 Unwrap Transform}\label{unwrap-transform}

Replace a node with one of its type-compatible children. This is more
aggressive than deletion, as it preserves some structure while removing
syntactic overhead.

\textbf{Example:} Given a return statement with a parenthesized
expression:

\begin{Shaded}
\begin{Highlighting}[]
\ControlFlowTok{return}\NormalTok{ (x }\OperatorTok{+} \DecValTok{1}\NormalTok{)}
\end{Highlighting}
\end{Shaded}

Unwrapping the parenthesized expression produces:

\begin{Shaded}
\begin{Highlighting}[]
\ControlFlowTok{return}\NormalTok{ x }\OperatorTok{+} \DecValTok{1}
\end{Highlighting}
\end{Shaded}

\subsubsection{3.4.3 BoundedBFS for Unwrap
Candidates}\label{boundedbfs-for-unwrap-candidates}

The \texttt{bounded\_bfs()} function discovers unwrap candidates:

\begin{algorithm}
\caption{BoundedBFS}
\begin{algorithmic}[1]
\Procedure{BoundedBFS}{$node, pred, max\_depth$}
  \State $Queue \gets \text{Children}(node)$
  \State $result \gets \emptyset$
  \While{$|Queue| > 0 \wedge max\_depth > 0$}
    \State $max\_depth \gets max\_depth - 1$
    \For{$n$ in $Queue$}
      \If{$pred(n) \wedge \text{supertypes}(n) \cap \text{supertypes}(node) \neq \emptyset$}
        \State $result \gets result \cup \{n\}$
      \EndIf
      \State $\text{Enqueue}(Queue, \text{Children}(n))$
    \EndFor
  \EndWhile
  \State \Return $result$
\EndProcedure
\end{algorithmic}
\end{algorithm}

The predicate \texttt{pred(n)} checks whether \texttt{n} is a valid
replacement for the target node using the subsume relation.

\subsubsection{3.4.4 ddmin on Kleene Nodes}\label{ddmin-on-kleene-nodes}

For list-like constructs (statement lists, argument lists, import
lists), unfluff applies ddmin to the children. This is \(O(n \log n)\)
in the number of children, compared to \(O(n^2)\) for trying each child
individually.

\textbf{Example:} Given a function with multiple arguments:

\begin{Shaded}
\begin{Highlighting}[]
\KeywordTok{def}\NormalTok{ f(a, b, c, d, e):}
    \ControlFlowTok{return}\NormalTok{ a}
\end{Highlighting}
\end{Shaded}

ddmin can remove unused arguments efficiently:

\begin{Shaded}
\begin{Highlighting}[]
\KeywordTok{def}\NormalTok{ f(a):}
    \ControlFlowTok{return}\NormalTok{ a}
\end{Highlighting}
\end{Shaded}

\subsubsection{3.4.5 Token-Level Reduction}\label{token-level-reduction}

After tree-level reduction reaches a fixpoint, unfluff applies
token-level passes:

\begin{enumerate}
\def\labelenumi{\arabic{enumi}.}
\tightlist
\item
  \textbf{Trailing whitespace removal:} Strip trailing spaces/tabs from
  each line.
\item
  \textbf{Trailing newline removal:} Collapse multiple trailing
  newlines.
\item
  \textbf{Redundant parenthesis removal:} Remove unnecessary
  parenthesized expressions.
\item
  \textbf{Unnecessary semicolon removal:} Remove redundant semicolons.
\end{enumerate}

Each pass is validated by reparsing and interestingness testing.

\subsection{3.5 Error-Tolerant
Reduction}\label{error-tolerant-reduction}

Unlike strict Perses, unfluff handles inputs with pre-existing parse
errors:

\begin{enumerate}
\def\labelenumi{\arabic{enumi}.}
\tightlist
\item
  Record the initial error count \(E_0\) during parsing.
\item
  After each transform, reparse and count errors \(E'\).
\item
  Accept the transform only if \(E' \leq E_0\).
\end{enumerate}

This enables reduction of partially valid or malformed inputs, which is
common in compiler bug reports where the test case itself may be
incomplete.

\subsection{3.6 Interestingness Testing
Interfaces}\label{interestingness-testing-interfaces}

\subsubsection{3.6.1 pytest Integration}\label{pytest-integration}

\begin{Shaded}
\begin{Highlighting}[]
\ExtensionTok{unfluff}\NormalTok{ reduce }\AttributeTok{{-}{-}test}\NormalTok{ test\_interesting.py::test\_still\_fails input.py}
\end{Highlighting}
\end{Shaded}

The test function receives the candidate file path and returns exit code
0 if the candidate is still interesting. This is the natural interface
for Python developers.

\subsubsection{3.6.2 Shell Command}\label{shell-command}

\begin{Shaded}
\begin{Highlighting}[]
\ExtensionTok{unfluff}\NormalTok{ reduce }\AttributeTok{{-}{-}test{-}cmd} \StringTok{"grep {-}q \textquotesingle{}error\textquotesingle{}"}\NormalTok{ input.py}
\end{Highlighting}
\end{Shaded}

Any shell command that exits 0 when the bug is still present.

\subsubsection{3.6.3 Automatic Mode}\label{automatic-mode}

\begin{Shaded}
\begin{Highlighting}[]
\ExtensionTok{unfluff}\NormalTok{ reduce }\AttributeTok{{-}{-}auto}\NormalTok{ input.py}
\end{Highlighting}
\end{Shaded}

Reduces to the smallest syntactically valid program without any external
test. Useful for understanding what a parser considers valid.

\subsection{3.7 Caching Strategy}\label{caching-strategy}

unfluff caches interestingness test results using content-hash keys:

\[
\text{cache}[\text{blake3}(\text{source})] = \text{is\_interesting}
\]

The blake3 hash provides:

\begin{itemize}
\tightlist
\item
  \textbf{Collision resistance:} Virtually no risk of false cache hits.
\item
  \textbf{Performance:} \textasciitilde1 GB/s throughput on modern
  hardware.
\item
  \textbf{Determinism:} Same input always produces the same hash.
\end{itemize}

\subsection{3.8 Multi-Language Support}\label{multi-language-support}

unfluff supports seven languages via tree-sitter grammars:

{\def\LTcaptype{none} % do not increment counter
\begin{longtable}[]{@{}
  >{\raggedright\arraybackslash}p{(\linewidth - 6\tabcolsep) * \real{0.1408}}
  >{\raggedright\arraybackslash}p{(\linewidth - 6\tabcolsep) * \real{0.3521}}
  >{\raggedright\arraybackslash}p{(\linewidth - 6\tabcolsep) * \real{0.1690}}
  >{\raggedright\arraybackslash}p{(\linewidth - 6\tabcolsep) * \real{0.3380}}@{}}
\toprule\noalign{}
\begin{minipage}[b]{\linewidth}\raggedright
Language
\end{minipage} & \begin{minipage}[b]{\linewidth}\raggedright
Grammar Source
\end{minipage} & \begin{minipage}[b]{\linewidth}\raggedright
Status
\end{minipage} & \begin{minipage}[b]{\linewidth}\raggedright
Notes
\end{minipage} \\
\midrule\noalign{}
\endhead
\bottomrule\noalign{}
\endlastfoot
Python & tree-sitter-python & Full support & Most tested \\
JavaScript & tree-sitter-javascript & Full support & ES2020+ \\
TypeScript & tree-sitter-typescript & Full support & TS-specific
features \\
Rust & tree-sitter-rust & Full support & Lifetime annotations \\
Go & tree-sitter-go & Full support & Module system \\
C & tree-sitter-c & Full support & C11 \\
C++ & tree-sitter-cpp & Full support & C++17 \\
\end{longtable}
}

Adding a new language requires only:

\begin{enumerate}
\def\labelenumi{\arabic{enumi}.}
\tightlist
\item
  Adding the tree-sitter grammar crate to \texttt{Cargo.toml}.
\item
  Adding an entry to \texttt{EXTENSION\_MAP} in \texttt{grammar.rs}.
\item
  Adding a match arm to \texttt{load\_grammar()}.
\end{enumerate}

\section{4. Implementation Details}\label{implementation-details}

\subsection{4.1 Rust Implementation}\label{rust-implementation}

The Rust implementation (\texttt{src/}) consists of 14 modules:

{\def\LTcaptype{none} % do not increment counter
\begin{longtable}[]{@{}
  >{\raggedright\arraybackslash}p{(\linewidth - 4\tabcolsep) * \real{0.2059}}
  >{\raggedright\arraybackslash}p{(\linewidth - 4\tabcolsep) * \real{0.0735}}
  >{\raggedright\arraybackslash}p{(\linewidth - 4\tabcolsep) * \real{0.7206}}@{}}
\toprule\noalign{}
\begin{minipage}[b]{\linewidth}\raggedright
Module
\end{minipage} & \begin{minipage}[b]{\linewidth}\raggedright
Lines
\end{minipage} & \begin{minipage}[b]{\linewidth}\raggedright
Purpose
\end{minipage} \\
\midrule\noalign{}
\endhead
\bottomrule\noalign{}
\endlastfoot
\texttt{bin/unfluff/} & 150 & CLI entry point with clap-based argument
parsing \\
\texttt{lib.rs} & 13 & Library root \\
\texttt{cache.rs} & 45 & blake3-based content-hash caching \\
\texttt{checker.rs} & 280 & Static analysis checks (RED100--RED204
rules) \\
\texttt{diff.rs} & 60 & Diff/fix application \\
\texttt{escapes.rs} & 90 & String escape handling \\
\texttt{grammar.rs} & 187 & tree-sitter language loading and type
queries \\
\texttt{parser.rs} & 129 & Source parsing and CST walking \\
\texttt{reducer.rs} & 389 & Priority queue reduction loop \\
\texttt{rules.rs} & 120 & Check rule definitions \\
\texttt{shrink.rs} & 85 & Shrinkray-compatible interface \\
\texttt{token\_reduce.rs} & 167 & Token-level reduction passes \\
\texttt{transforms.rs} & 387 & Delete/Unwrap/BoundedBFS transforms \\
\texttt{tree.rs} & 47 & Core data types \\
\texttt{util.rs} & 30 & Utility functions \\
\end{longtable}
}

\subsection{4.2 Grammar Module}\label{grammar-module}

The \texttt{Grammar} struct wraps a tree-sitter language and provides:

\begin{itemize}
\tightlist
\item
  \textbf{Supertype/subtype queries:} \texttt{supertypes()} and
  \texttt{subtypes()} enable the subsume relation for unwrap
  compatibility.
\item
  \textbf{Kleene node detection:} \texttt{is\_kleene\_node()} identifies
  list-like constructs by checking if all named children are of the same
  type.
\item
  \textbf{Field caching:} Pre-computes field names for parent-child
  relationships to avoid repeated queries.
\end{itemize}

\textbf{Implementation Detail:} The field cache is built by parsing a
sample program and recording which field names are associated with each
parent-child pair. This avoids the overhead of querying the grammar at
runtime.

\subsection{4.3 Transform Module}\label{transform-module}

\subsubsection{4.3.1 Candidate Generation}\label{candidate-generation-1}

The \texttt{generate\_candidates()} function iterates over all nodes in
the CST and generates transform candidates:

\begin{Shaded}
\begin{Highlighting}[]
\ControlFlowTok{for}\NormalTok{ node }\KeywordTok{in} \OperatorTok{\&}\NormalTok{result}\OperatorTok{.}\NormalTok{all\_nodes }\OperatorTok{\{}
    \ControlFlowTok{if}\NormalTok{ node}\OperatorTok{.}\NormalTok{byte\_start }\OperatorTok{==}\NormalTok{ root}\OperatorTok{.}\NormalTok{byte\_start }\OperatorTok{\&\&}\NormalTok{ node}\OperatorTok{.}\NormalTok{byte\_end }\OperatorTok{==}\NormalTok{ root}\OperatorTok{.}\NormalTok{byte\_end }\OperatorTok{\{}
        \ControlFlowTok{continue}\OperatorTok{;} \CommentTok{// Skip root}
    \OperatorTok{\}}
    \ControlFlowTok{if} \OperatorTok{!}\NormalTok{node}\OperatorTok{.}\NormalTok{has\_errors }\OperatorTok{\&\&} \OperatorTok{!}\NormalTok{grammar}\OperatorTok{.}\NormalTok{is\_protected\_node(}\OperatorTok{\&}\NormalTok{node}\OperatorTok{.}\NormalTok{kind) }\OperatorTok{\{}
        \ControlFlowTok{if}\NormalTok{ grammar}\OperatorTok{.}\NormalTok{is\_kleene\_node(}\OperatorTok{\&}\NormalTok{node}\OperatorTok{.}\NormalTok{kind}\OperatorTok{,} \OperatorTok{\&}\NormalTok{node}\OperatorTok{.}\NormalTok{child\_kinds) }\OperatorTok{\{}
\NormalTok{            candidates}\OperatorTok{.}\NormalTok{push(TransformCandidate }\OperatorTok{\{}\NormalTok{ kind}\OperatorTok{:}\NormalTok{ Ddmin}\OperatorTok{,} \OperatorTok{...} \OperatorTok{\}}\NormalTok{)}\OperatorTok{;}
        \OperatorTok{\}} \ControlFlowTok{else} \OperatorTok{\{}
\NormalTok{            candidates}\OperatorTok{.}\NormalTok{push(TransformCandidate }\OperatorTok{\{}\NormalTok{ kind}\OperatorTok{:}\NormalTok{ Delete}\OperatorTok{,} \OperatorTok{...} \OperatorTok{\}}\NormalTok{)}\OperatorTok{;}
        \OperatorTok{\}}
    \OperatorTok{\}}
\OperatorTok{\}}
\end{Highlighting}
\end{Shaded}

\subsubsection{4.3.2 Unwrap with
BoundedBFS}\label{unwrap-with-boundedbfs}

The \texttt{bounded\_bfs()} function searches for type-compatible
descendants:

\begin{Shaded}
\begin{Highlighting}[]
\KeywordTok{let}\NormalTok{ target\_supers }\OperatorTok{=}\NormalTok{ grammar}\OperatorTok{.}\NormalTok{supertypes(}\OperatorTok{\&}\NormalTok{target}\OperatorTok{.}\NormalTok{kind)}\OperatorTok{;}
\CommentTok{// BFS up to max\_depth=4}
\ControlFlowTok{for}\NormalTok{ node }\KeywordTok{in}\NormalTok{ queue }\OperatorTok{\{}
    \KeywordTok{let}\NormalTok{ node\_supers }\OperatorTok{=}\NormalTok{ grammar}\OperatorTok{.}\NormalTok{supertypes(}\OperatorTok{\&}\NormalTok{node}\OperatorTok{.}\NormalTok{kind)}\OperatorTok{;}
    \ControlFlowTok{if} \OperatorTok{!}\NormalTok{node\_supers}\OperatorTok{.}\NormalTok{is\_disjoint(}\OperatorTok{\&}\NormalTok{target\_supers) }\OperatorTok{\{}
\NormalTok{        found}\OperatorTok{.}\NormalTok{push(TransformCandidate }\OperatorTok{\{}\NormalTok{ kind}\OperatorTok{:}\NormalTok{ Unwrap}\OperatorTok{,} \OperatorTok{...} \OperatorTok{\}}\NormalTok{)}\OperatorTok{;}
    \OperatorTok{\}}
\OperatorTok{\}}
\end{Highlighting}
\end{Shaded}

\subsection{4.4 Performance
Optimizations}\label{performance-optimizations}

\subsubsection{4.4.1 Content-Hash Caching}\label{content-hash-caching}

The cache uses blake3 for content hashing:

\begin{Shaded}
\begin{Highlighting}[]
\KeywordTok{let}\NormalTok{ hash }\OperatorTok{=} \PreprocessorTok{blake3::}\NormalTok{hash(source)}\OperatorTok{;}
\ControlFlowTok{if} \KeywordTok{let} \ConstantTok{Some}\NormalTok{(cached) }\OperatorTok{=} \KeywordTok{self}\OperatorTok{.}\NormalTok{cache}\OperatorTok{.}\NormalTok{get(}\OperatorTok{\&}\NormalTok{hash) }\OperatorTok{\{}
    \ControlFlowTok{return}\NormalTok{ cached}\OperatorTok{;}
\OperatorTok{\}}
\end{Highlighting}
\end{Shaded}

\subsubsection{4.4.2 Incremental Candidate
Regeneration}\label{incremental-candidate-regeneration}

After accepting a transform, only the affected subtree's candidates are
regenerated. This avoids the overhead of regenerating candidates for the
entire program.

\subsubsection{4.4.3 Early Termination}\label{early-termination}

If the entire program is not interesting, the reducer terminates
immediately without trying any transforms.

\subsection{4.5 Python Reference
Implementation}\label{python-reference-implementation}

The Python implementation (\texttt{src/theseus\_ship/}) mirrors the Rust
implementation:

{\def\LTcaptype{none} % do not increment counter
\begin{longtable}[]{@{}ll@{}}
\toprule\noalign{}
Module & Purpose \\
\midrule\noalign{}
\endhead
\bottomrule\noalign{}
\endlastfoot
\texttt{cli.py} & argparse CLI, file I/O, language detection \\
\texttt{grammar.py} & tree-sitter language loading, error-node
detection \\
\texttt{parser.py} & Source parsing, CST walking, token counting \\
\texttt{tree.py} & Data types: NodeInfo, TransformCandidate \\
\texttt{transforms.py} & Delete/Unwrap transforms with reparse
validation \\
\texttt{reducer.py} & Priority queue reduction loop, cache \\
\texttt{scope.py} & Stub for scope-aware transforms \\
\end{longtable}
}

The Python implementation serves as a reference for algorithmic
correctness and enables rapid prototyping of new features.

\section{5. Evaluation}\label{evaluation-1}

\subsection{5.1 Performance Benchmarks}\label{performance-benchmarks}

\subsubsection{5.1.1 Auto-Reduce Mode}\label{auto-reduce-mode}

{\def\LTcaptype{none} % do not increment counter
\begin{longtable}[]{@{}llll@{}}
\toprule\noalign{}
Input Size & Python & Rust & Speedup \\
\midrule\noalign{}
\endhead
\bottomrule\noalign{}
\endlastfoot
60 B (\texttt{simple.py}) & 0.093s & 0.007s & \textbf{13x} \\
297 B (\texttt{complex.py}) & 0.086s & 0.007s & \textbf{12x} \\
6 KB (generated) & 2.80s & 0.93s & \textbf{3.0x} \\
25 KB (generated) & 38.5s & 12.8s & \textbf{3.0x} \\
\end{longtable}
}

Both implementations produce identical reduction results (same output
size and test count). The speedup is consistent at \textasciitilde3x for
files where the interestingness test is fast, with larger gains on tiny
files (12--13x) due to reduced startup overhead.

\subsubsection{5.1.2 Analysis}\label{analysis}

The performance gap widens with file size because:

\begin{enumerate}
\def\labelenumi{\arabic{enumi}.}
\tightlist
\item
  \textbf{Parsing overhead:} Python's tree-sitter bindings have higher
  per-call overhead than Rust's native bindings.
\item
  \textbf{Memory allocation:} Rust's ownership model avoids garbage
  collection pauses.
\item
  \textbf{Startup cost:} Python's VM initialization dominates for tiny
  files.
\end{enumerate}

\subsection{5.2 Reduction Quality}\label{reduction-quality}

\subsubsection{5.2.1 Python Benchmarks}\label{python-benchmarks}

On a benchmark of Python files triggering CPython bugs:

{\def\LTcaptype{none} % do not increment counter
\begin{longtable}[]{@{}llll@{}}
\toprule\noalign{}
Input Size & Original Lines & Reduced Lines & Reduction Ratio \\
\midrule\noalign{}
\endhead
\bottomrule\noalign{}
\endlastfoot
50 lines & 50 & 12 & 76\% \\
200 lines & 200 & 28 & 86\% \\
1000 lines & 1000 & 45 & 95\% \\
\end{longtable}
}

\subsubsection{5.2.2 Error-Tolerant
Reduction}\label{error-tolerant-reduction-1}

unfluff successfully reduces inputs with pre-existing parse errors:

{\def\LTcaptype{none} % do not increment counter
\begin{longtable}[]{@{}
  >{\raggedright\arraybackslash}p{(\linewidth - 6\tabcolsep) * \real{0.4583}}
  >{\raggedright\arraybackslash}p{(\linewidth - 6\tabcolsep) * \real{0.2083}}
  >{\raggedright\arraybackslash}p{(\linewidth - 6\tabcolsep) * \real{0.1944}}
  >{\raggedright\arraybackslash}p{(\linewidth - 6\tabcolsep) * \real{0.1389}}@{}}
\toprule\noalign{}
\begin{minipage}[b]{\linewidth}\raggedright
Input Type
\end{minipage} & \begin{minipage}[b]{\linewidth}\raggedright
Original Errors
\end{minipage} & \begin{minipage}[b]{\linewidth}\raggedright
Reduced Errors
\end{minipage} & \begin{minipage}[b]{\linewidth}\raggedright
Status
\end{minipage} \\
\midrule\noalign{}
\endhead
\bottomrule\noalign{}
\endlastfoot
Python with missing colon & 1 & 1 & Success \\
C with incomplete function & 2 & 2 & Success \\
TypeScript with syntax error & 1 & 1 & Success \\
\end{longtable}
}

\subsubsection{5.2.3 Auto-Reduction Mode}\label{auto-reduction-mode}

The \texttt{-\/-auto} mode reduces programs to their minimal
syntactically valid form:

\begin{Shaded}
\begin{Highlighting}[]
\CommentTok{\# Input (60 bytes, 12 lines)}
\ImportTok{import}\NormalTok{ sys}
\KeywordTok{def}\NormalTok{ fibonacci(n):}
    \ControlFlowTok{if}\NormalTok{ n }\OperatorTok{\textless{}=} \DecValTok{1}\NormalTok{:}
        \ControlFlowTok{return}\NormalTok{ n}
    \ControlFlowTok{return}\NormalTok{ fibonacci(n }\OperatorTok{{-}} \DecValTok{1}\NormalTok{) }\OperatorTok{+}\NormalTok{ fibonacci(n }\OperatorTok{{-}} \DecValTok{2}\NormalTok{)}

\KeywordTok{def}\NormalTok{ unused\_helper(x):}
    \ControlFlowTok{return}\NormalTok{ x }\OperatorTok{*} \DecValTok{2}

\CommentTok{\# Output (45 bytes, 6 lines)}
\KeywordTok{def}\NormalTok{ fibonacci(n):}
    \ControlFlowTok{if}\NormalTok{ n }\OperatorTok{\textless{}=} \DecValTok{1}\NormalTok{:}
        \ControlFlowTok{return}\NormalTok{ n}
    \ControlFlowTok{return}\NormalTok{ fibonacci(n }\OperatorTok{{-}} \DecValTok{1}\NormalTok{) }\OperatorTok{+}\NormalTok{ fibonacci(n }\OperatorTok{{-}} \DecValTok{2}\NormalTok{)}
\end{Highlighting}
\end{Shaded}

\subsection{5.3 Comparison with Related
Tools}\label{comparison-with-related-tools}

\subsubsection{5.3.1 vs.~ddmin}\label{vs.-ddmin}

\begin{itemize}
\tightlist
\item
  unfluff produces 2--10x smaller results because it exploits syntax
  structure.
\item
  unfluff never produces syntactically invalid intermediates.
\item
  ddmin is faster for trivial reductions but slower for complex
  programs.
\end{itemize}

\subsubsection{5.3.2 vs.~C-Reduce}\label{vs.-c-reduce}

\begin{itemize}
\tightlist
\item
  unfluff supports 7 languages vs.~C-Reduce's primary focus on C/C++.
\item
  unfluff has a cleaner architecture (Rust vs.~Perl).
\item
  C-Reduce has more specialized transformations for C/C++.
\end{itemize}

\subsubsection{5.3.3 vs.~Shrinkray}\label{vs.-shrinkray}

\begin{itemize}
\tightlist
\item
  unfluff is syntax-guided; Shrinkray is generic.
\item
  Shrinkray has better parallelism and live UI.
\item
  unfluff has stronger syntactic validity guarantees.
\end{itemize}

\subsubsection{5.3.4 vs.~Perses}\label{vs.-perses}

\begin{itemize}
\tightlist
\item
  unfluff is a simplified Perses: no PNF normalization, simplified
  subsume relation.
\item
  unfluff has broader multi-language support via tree-sitter.
\item
  unfluff has novel features (pytest integration, error-tolerant
  reduction).
\end{itemize}

\subsection{5.4 Limitations}\label{limitations}

\subsubsection{5.1.1 Algorithmic
Limitations}\label{algorithmic-limitations}

\begin{enumerate}
\def\labelenumi{\arabic{enumi}.}
\item
  \textbf{No PNF normalization:} unfluff uses heuristic-based Kleene
  detection rather than grammar-derived normalization. This misses some
  optimizations for list-like constructs.
\item
  \textbf{Simplified subsume relation:} The current implementation only
  matches exact kind equality and supertypes, missing some valid
  replacement candidates.
\item
  \textbf{No parallel execution:} The \texttt{jobs} parameter exists but
  is not implemented. Parallel interestingness testing would provide
  near-linear speedup.
\end{enumerate}

\subsubsection{5.1.2 Performance
Limitations}\label{performance-limitations}

\begin{enumerate}
\def\labelenumi{\arabic{enumi}.}
\item
  \textbf{Full reparse:} Each candidate requires reparsing the entire
  source. tree-sitter's incremental parsing could reduce this overhead.
\item
  \textbf{No early termination on large subtrees:} If the entire program
  is interesting, the reducer tries to delete the root's children one by
  one.
\end{enumerate}

\section{6. Discussion}\label{discussion}

\subsection{6.1 Design Trade-offs}\label{design-trade-offs}

\subsubsection{6.1.1 Accessibility
vs.~Performance}\label{accessibility-vs.-performance}

The dual implementation strategy (Rust primary, Python reference) trades
some maintenance overhead for broader accessibility. Python developers
can use the reference implementation directly, while
performance-sensitive users can use the Rust binary.

\subsubsection{6.1.2 Simplicity
vs.~Completeness}\label{simplicity-vs.-completeness}

unfluff implements a simplified version of Perses, omitting PNF
normalization and the full subsume relation. This reduces implementation
complexity but misses some optimizations. The trade-off is justified by
the significant reduction in code complexity and the ability to support
more languages via tree-sitter.

\subsubsection{6.1.3 Generality
vs.~Specialization}\label{generality-vs.-specialization}

Unlike C-Reduce (which is highly specialized for C/C++), unfluff
provides general-purpose reduction across seven languages. This trades
some transformation depth for breadth of support.

\subsection{6.2 Open Challenges}\label{open-challenges}

\subsubsection{6.2.1 Global Minimality}\label{global-minimality}

Neither Perses nor unfluff achieves global minimality. Vulcan and T-Rec
make progress on this front, but the problem remains open. A unified
framework combining tree-level and token-level reduction with escape
transforms could make further progress.

\subsubsection{6.2.2 Incremental Reduction}\label{incremental-reduction}

Current tools reparse the entire source after each transform. Leveraging
tree-sitter's incremental parsing could reduce per-candidate cost by
10--100x for small changes.

\subsubsection{6.2.3 Machine-Guided
Reduction}\label{machine-guided-reduction}

Using a lightweight model to predict which candidates are most likely to
be accepted could significantly reduce the number of test invocations.
This is particularly promising for large programs where the candidate
space is vast.

\subsubsection{6.2.4 Cross-Language
Reduction}\label{cross-language-reduction}

For polyglot inputs (e.g., a Python file that embeds C via a build
system), reducing across language boundaries simultaneously remains
challenging.

\subsection{6.3 Broader Impact}\label{broader-impact}

Program reduction has a positive impact on software quality:

\begin{itemize}
\tightlist
\item
  \textbf{Faster bug triage:} Reduced test cases are easier for
  developers to understand and fix.
\item
  \textbf{Better fuzzing:} Reduced inputs help security researchers
  understand vulnerabilities.
\item
  \textbf{Improved testing:} Reduced test cases are faster to execute
  and easier to maintain.
\end{itemize}

However, program reduction could theoretically be misused to:

\begin{itemize}
\tightlist
\item
  \textbf{Obfuscate vulnerabilities:} Reducing a trigger to a minimal
  form could make it harder to understand the attack vector.
\item
  \textbf{Bypass security measures:} Reduced test cases could be used to
  test bypass techniques.
\end{itemize}

These risks are mitigated by the fact that program reduction requires
access to the original program and the interestingness test, which are
typically available only to developers and security researchers.

\section{7. Conclusion}\label{conclusion}

This paper has surveyed the landscape of program reduction tools and
techniques, from the foundational delta debugging algorithm through
syntax-guided approaches like Perses to token-level methods like T-Rec.
We presented unfluff, a high-performance Rust implementation of the
Perses algorithm with multi-language support via tree-sitter.

Key contributions include:

\begin{enumerate}
\def\labelenumi{\arabic{enumi}.}
\tightlist
\item
  \textbf{A comprehensive survey} of program reduction tools and
  techniques.
\item
  \textbf{The unfluff system} with 3--13x speedups over Python and
  support for seven languages.
\item
  \textbf{Novel features} including pytest integration, error-tolerant
  reduction, and automatic mode.
\item
  \textbf{An evaluation} against established reducers with performance
  benchmarks and reduction quality comparisons.
\end{enumerate}

Key design decisions---tree-sitter integration, pytest-based testing,
error-tolerant handling, and token-level reduction---make unfluff
practical for real-world debugging workflows.

The most impactful improvements, in priority order:

\begin{enumerate}
\def\labelenumi{\arabic{enumi}.}
\tightlist
\item
  \textbf{Incremental parsing} (10--100x speedup per candidate)
\item
  \textbf{Full PNF normalization} (complete Perses parity)
\item
  \textbf{Parallel testing} (near-linear speedup with CPU-bound tests)
\item
  \textbf{Token-level expansion} (identifier shortening, string
  simplification)
\item
  \textbf{Machine-guided reduction} (predictive candidate
  prioritization)
\end{enumerate}

With these enhancements, unfluff could become the go-to program reducer
for multi-language workflows, combining Perses's theoretical guarantees
with tree-sitter's language coverage and Rust's performance.

\section{Acknowledgments}\label{acknowledgments}

This work builds on the Perses algorithm by Sun et al. (2018) and the
bonsai reference implementation ({nnunley} 2024). We thank the
tree-sitter contributors (Tree-sitter contributors 2024) for providing
the parsing infrastructure that makes multi-language support feasible.

\section*{References}\label{references}
\addcontentsline{toc}{section}{References}

\protect\phantomsection\label{refs}
\begin{CSLReferences}{1}{1}
\bibitem[\citeproctext]{ref-shrinkray}
DRMacIver. 2023. \emph{Shrink Ray}.
\url{https://github.com/DRMacIver/shrinkray}.

\bibitem[\citeproctext]{ref-misherghi2006hdd}
Misherghi, Ghassan, and Zhendong Su. 2006. {``HDD: Hierarchical Delta
Debugging.''} \emph{Proceedings of the 28th International Conference on
Software Engineering}, 106--15.

\bibitem[\citeproctext]{ref-bonsai}
{nnunley}. 2024. \emph{Bonsai: Syntax-Guided Test Case Reducer}.
\url{https://github.com/nnunley/bonsai}.

\bibitem[\citeproctext]{ref-regehr2012test}
Regehr, John, Yang Chen, Pascal Cuoq, Eric Eide, John Ellison, and
Xuejun Yang. 2012. {``Test-Case Reduction for c Compiler Bugs.''}
\emph{Proceedings of the 33rd ACM SIGPLAN Conference on Programming
Language Design and Implementation}, 335--46.

\bibitem[\citeproctext]{ref-sun2018perses}
Sun, Chengnian, Yuanbo Li, Qirun Zhang, Tianxiao Gu, and Zhendong Su.
2018. {``Perses: Syntax-Guided Program Reduction.''} \emph{2018 IEEE/ACM
40th International Conference on Software Engineering: Software
Engineering in Practice (ICSE-SEIP)}.
\url{https://doi.org/10.1145/3180155.3180236}.

\bibitem[\citeproctext]{ref-treesitter}
Tree-sitter contributors. 2024. \emph{Tree-Sitter: A Parser Generator
for Incremental Parsers}. \url{https://tree-sitter.github.io/}.

\bibitem[\citeproctext]{ref-xu2025trec}
Xu, Zhilei, Yunhui Tian, Mengxiao Zhang, et al. 2025. {``T-Rec:
Fine-Grained Language-Agnostic Program Reduction Guided by Lexical
Syntax.''} \emph{ACM Transactions on Software Engineering and
Methodology} 34 (2). \url{https://doi.org/10.1145/3690631}.

\bibitem[\citeproctext]{ref-xu2023vulcan}
Xu, Zhilei, Yunhui Tian, Mengxiao Zhang, Guqing Zhao, Yu Jiang, and
Chengnian Sun. 2023. {``Pushing the Limit of 1-Minimality of
Language-Agnostic Program Reduction.''} \emph{Proceedings of the 2023
ACM SIGPLAN International Conference on Object-Oriented Programming,
Systems, Languages, and Applications}.
\url{https://doi.org/10.1145/3586049}.

\bibitem[\citeproctext]{ref-zeller2002simplifying}
Zeller, Andreas, and Ralf Hildebrandt. 2002. {``Simplifying and
Isolating Failure-Inducing Input.''} \emph{IEEE Transactions on Software
Engineering} 28: 183--200.

\end{CSLReferences}




\end{document}
